\textbf{Phần 1. Trắc nghiệm lựa chọn}
\setcounterpageref{ex}{0}

\begin{ex} %[2D14N71]
Kết quả của phép tính $2 \times 7$ là
\choice
{$12$}
{$16$}
{\True $14$}
{$18$}
\loigiai{
$2 \times 7 = 14$.\\
\textbf{Đáp án của câu hỏi là C}
}
\end{ex}

\begin{ex} %[2D14N71]
Kết quả của phép tính $5 \times 8$ là
\choice
{$35$}
{$45$}
{$50$}
{\True $40$}
\loigiai{
$5 \times 8 = 40$.\\
\textbf{Đáp án của câu hỏi là D}
}
\end{ex}

\begin{ex} %[2D14H71]
Kết quả của phép tính $18 : 2$ là
\choice
{\True $9$}
{$8$}
{$7$}
{$6$}
\loigiai{
$18 : 2 = 9$.\\
\textbf{Đáp án của câu hỏi là A}
}
\end{ex}

\begin{ex} %[2D14H71]
Kết quả của phép tính $45 : 5$ là
\choice
{$8$}
{\True $9$}
{$7$}
{$6$}
\loigiai{
$45 : 5 = 9$.\\
\textbf{Đáp án của câu hỏi là B}
}
\end{ex}

\begin{ex} %[2D14H71]
Trong phép nhân $5 \times 9 = 45$, số 45 được gọi là
\choice
{Thừa số}
{Số bị chia}
{Thương}
{\True Tích}
\loigiai{
45 là Tích.\\
\textbf{Đáp án của câu hỏi là D}
}
\end{ex}

\begin{ex} %[2D14H71]
Trong phép chia $20 : 2 = 10$, số 2 được gọi là
\choice
{Số bị chia}
{\True Số chia}
{Thương}
{Tích}
\loigiai{
2 là Số chia.\\
\textbf{Đáp án của câu hỏi là B}
}
\end{ex}

\begin{ex} %[2D14V71]
Mỗi bạn có 5 cái nhãn vở. Hỏi 7 bạn có tất cả bao nhiêu cái nhãn vở?
\choice
{\True $35$ cái}
{$30$ cái}
{$40$ cái}
{$45$ cái}
\loigiai{
$5 \times 7 = 35$ cái.\\
\textbf{Đáp án của câu hỏi là A}
}
\end{ex}

\begin{ex} %[2D14V71]
Có 16 cái bánh chia đều cho 2 bạn. Mỗi bạn được số cái bánh là
\choice
{$7$ cái}
{$9$ cái}
{\True $8$ cái}
{$10$ cái}
\loigiai{
$16 : 2 = 8$ cái.\\
\textbf{Đáp án của câu hỏi là C}
}
\end{ex}

\begin{ex} %[2D14V71]
Có 35 lít nước mắm rót đều vào 5 can. Mỗi can có số lít nước mắm là
\choice
{$6\text{ l}$}
{$8\text{ l}$}
{$9\text{ l}$}
{\True $7\text{ l}$}
\loigiai{
$35 : 5 = 7\text{ l}$.\\
\textbf{Đáp án của câu hỏi là D}
}
\end{ex}

\begin{ex} %[2D14V71]
Tính giá trị biểu thức: $2 \times 9 + 32$.
\choice
{\True $50$}
{$48$}
{$52$}
{$54$}
\loigiai{
$18 + 32 = 50$.\\
\textbf{Đáp án của câu hỏi là A}
}
\end{ex}

\begin{ex} %[2D14V71]
Tính giá trị biểu thức: $50 : 5 - 4$.
\choice
{$5$}
{\True $6$}
{$7$}
{$8$}
\loigiai{
$10 - 4 = 6$.\\
\textbf{Đáp án của câu hỏi là B}
}
\end{ex}

\begin{ex} %[2D14C71]
Một lớp có 30 học sinh chia đều thành 5 tổ. Mỗi tổ có 2 bạn nữ. Hỏi cả lớp có bao nhiêu bạn nam?
\choice
{$18$ bạn}
{$22$ bạn}
{\True $20$ bạn}
{$24$ bạn}
\loigiai{
Số bạn nữ cả lớp: $2 \times 5 = 10$ bạn.\\
Số bạn nam cả lớp: $30 - 10 = 20$ bạn.\\
\textbf{Đáp án của câu hỏi là C}
}
\end{ex}

\textbf{Phần 2. Trắc nghiệm đúng sai}
\setcounterpageref{ex}{0}

\begin{ex} %[2D14H71]
Xét tính đúng sai của các phép toán nhân chia:
\choiceTF
{\True $2 \times 8 = 16$.}
{\True $5 \times 6 = 30$.}
{$14 : 2 = 8$.}
{\True $14 : 2 = 7$.}
\loigiai{
\begin{itemize}
    \item $2 \times 8 = 16$: Đúng.
    \item $5 \times 6 = 30$: Đúng.
    \item $14 : 2 = 7 \neq 8$. Ý 3 Sai, Ý 4 Đúng.
\end{itemize}
\begin{itemize}
    \item a) \textbf{Đúng}.
    \item b) \textbf{Đúng}.
    \item c) \textbf{Sai}.
    \item d) \textbf{Đúng}.
\end{itemize}
}
\end{ex}

\begin{ex} %[2D14V71]
Có 40 cái kẹo chia đều vào 5 hộp. Xét tính đúng sai:
\choiceTF
{\True Mỗi hộp có 8 cái kẹo ($40 : 5 = 8$).}
{\True 3 hộp như thế có 24 cái kẹo ($8 \times 3 = 24$).}
{Nếu bớt đi 1 hộp thì còn lại 30 cái kẹo.}
{\True Khi bớt 1 hộp thì còn lại $8 \times 4 = 32$ cái kẹo.}
\loigiai{
\begin{itemize}
    \item Mỗi hộp 8 cái: Đúng.
    \item 3 hộp có 24 cái: Đúng.
    \item $8 \times 4 = 32 \neq 30$. Ý 3 Sai, Ý 4 Đúng.
\end{itemize}
\begin{itemize}
    \item a) \textbf{Đúng}.
    \item b) \textbf{Đúng}.
    \item c) \textbf{Sai}.
    \item d) \textbf{Đúng}.
\end{itemize}
}
\end{ex}

\textbf{Phần 3. Trả lời ngắn}
\setcounterpageref{ex}{0}

\begin{ex} %[2D14H71]
Tính: $2 \times 9$.
\shortans{18}
\loigiai{
\textbf{Đáp số:} $18$\\
$2 \times 9 = 18$.
}
\end{ex}

\begin{ex} %[2D14H71]
Tính: $5 \times 9$.
\shortans{45}
\loigiai{
\textbf{Đáp số:} $45$\\
$5 \times 9 = 45$.
}
\end{ex}

\begin{ex} %[2D14V71]
Tính: $18 : 2$.
\shortans{9}
\loigiai{
\textbf{Đáp số:} $9$\\
$18 : 2 = 9$.
}
\end{ex}

\begin{ex} %[2D14V71]
Tính: $40 : 5$.
\shortans{8}
\loigiai{
\textbf{Đáp số:} $8$\\
$40 : 5 = 8$.
}
\end{ex}

\begin{ex} %[2D14V71]
Tính giá trị biểu thức: $5 \times 7 + 25$.
\shortans{60}
\loigiai{
\textbf{Đáp số:} $60$\\
$35 + 25 = 60$.
}
\end{ex}

\begin{ex} %[2D14C71]
Có 5 túi bi, mỗi túi 4 viên bi. Sau khi cho bạn 6 viên bi thì còn lại bao nhiêu viên bi?
\shortans{14}
\loigiai{
\textbf{Đáp số:} $14$\\
$5 \times 4 - 6 = 20 - 6 = 14$ viên bi.
}
\end{ex}
